

\section{Admissibility} \label{section:Admissibility}

In this section, we discuss the relation between admissibility and discrete decomposability in the branching problem.
We show in Theorem \ref{thm:DiscretelyDecomposableAndAdmissibility} that the two notions are equivalent modulo compact factors.
We then use the insight obtained from Theorem \ref{thm:DiscretelyDecomposableAndAdmissibility} to give an alternative proof of Fact \ref{fact:Implication} (1) $\Rightarrow$ (2).

\subsection{Criteria for admissibility}

Let $G$ be a connected reductive Lie group with a Cartan involution $\theta$ and $H$ a $\theta$-stable closed connected reductive subgroup of $G$.
Set $K\coloneq G^\theta$ and $K_H \coloneq K \cap H$.

We shall recall results on admissibility in the branching problem.
The `only if' part of the following fact was proved by T.\ Kobayashi in \cite[Proposition 1.6]{Ko98_discrete_decomposable_3}.
The `if' part was proved by Zhu--Liang in \cite{ZhLi10} for unitarizable irreducible modules, and by the author in \cite{Ki24} for general irreducible modules.

\begin{fact} \label{fact:AdmissibleKH}
	Let $V$ be an irreducible $(\lie{g}, K)$-module.
	Then $V$ is $K_H$-admissible if and only if $V$ is $\lie{h}$-admissible.
\end{fact}

The following theorem was proved in \cite[Theorem 20]{Ki24}.
For compact subgroups, the theorem was proved by T.\ Kobayashi in \cite[Theorem 2.9]{Ko98_discrete_decomposable_2} and \cite[Theorem 1.1]{Ko19_admissible}.

\begin{theorem} \label{thm:AdmissibilityCriterion}
	Let $V$ be an irreducible $(\lie{g}, K)$-module.
	Then $V$ is $\lie{h}$-admissible if and only if $\lie{h}^\perp \cap \WF(V) = \set{0}$.
\end{theorem}

Theorem \ref{thm:AdmissibilityCriterion} asserts that admissibility is equivalent to a purely Lie-algebraic condition on nilpotent orbits.
We shall show a nilpotent-orbit counterpart of Fact \ref{fact:AdmissibleKH}.

\begin{proposition} \label{prop:AdmissibilityCriterionWavefront}
	Let $\orbit{O}$ be a nilpotent $G$-orbit in $\lie{g}$.
	Then $\overline{\orbit{O}} \cap \lie{h}^\perp = \set{0}$ if and only if $\overline{\orbit{O}} \cap \lie{k}_H^\perp = \set{0}$.
\end{proposition}

\begin{proof}
	The `if' part is obvious from $\lie{k}_H \subset \lie{h}$.
	To show the converse, we assume that $\overline{\orbit{O}} \cap \lie{k}_H^\perp \neq \set{0}$.
	Note that $\lie{k}_H^\perp = \lie{h}^{-\theta} \oplus \lie{h}^\perp$.
	Take a non-zero $X\coloneq X_1 + X_2 \in \overline{\orbit{O}} \cap \lie{k}_H^\perp$ with $X_1 \in \lie{h}^{-\theta}$ and $X_2 \in \lie{h}^\perp$.
	
	Observe that $\ad(X_1)$ is semisimple and has real eigenvalues.
	Since $\overline{\orbit{O}}$ and $\lie{h}^\perp$ are $\exp(\RR \ad(X_1))$-stable, the same argument as in the proof of Proposition \ref{prop:NullCone} allows us to assume that $X$ is an eigenvector of $\ad(X_1)$, that is, $X$ has one of the following forms:
	\begin{enumparen}
		\item $X = X_1$,
		\item $X = X_2$,
		\item $X_1, X_2 \neq 0$ and $[X_1, X_2] = 0$.
	\end{enumparen}
	If $X = X_1$, then $X$ is semisimple, which contradicts the assumption that $X$ is non-zero and nilpotent.
	If $X = X_2$, then $0 \neq X \in \overline{\orbit{O}} \cap \lie{h}^\perp$ and we have nothing to prove.

	Assume that $X_1, X_2 \neq 0$ and $[X_1, X_2] = 0$.
	Since $X_1 \in \lie{h}$ and $X_2 \in \lie{h}^\perp$, we have $(X_1, X_2) = 0$.
	Since $X$ is a nilpotent element commuting with the semisimple element $X_1$, they are orthogonal.
	In fact, $X$ belongs to the derived subalgebra of $\lieCent_{\lie{g}}(X_1)$.
	Hence we have
	\begin{align*}
		0 = (X_1, X_2) = (X_1, X - X_1) = -(X_1, X_1).
	\end{align*}
	Since $X_1 \in \lie{h}^{-\theta}$, we obtain $X_1 = 0$, which is a contradiction.
	We have shown the proposition.
\end{proof}

Let $\lie{h}_\nc$ be the ideal of $\lie{h}$ generated by $\lie{h}^{-\theta}$ and $H_\nc$ the analytic subgroup of $H$ with the Lie algebra $\lie{h}_\nc$.
In other words, $\lie{h}_{\nc}$ is the non-compact part of $\lie{h}$.
Then $H_\nc$ is a closed normal subgroup of $H$ and $H/H_\nc$ is compact.
Set $H' \coloneq N_{G}(H_{\nc})_o$ and $K_{H'} \coloneq K \cap H'$.
Then we have
\begin{align*}
	\lie{h}' &= \lieNorm_{\lie{g}}(\lie{h}_\nc) = \lie{h}_\nc + \lieCent_{\lie{g}}(\lie{h}_\nc) = \lie{h}_\nc + \lie{g}^{H_{\nc}}, \\
	\lie{h}_{\nc}^\perp &= (\lie{h'})^\perp \oplus (\lie{h}_\nc^\perp \cap (\lie{h}')^{H_{\nc}}), \\
	\lie{h}_{\nc}^\perp &= \lie{h}^\perp \oplus (\lie{h}_\nc^\perp \cap \lie{h}^{H_{\nc}}).
\end{align*}

\begin{lemma} \label{lem:NilpotentAndNonCompact}
	One has
	\begin{align*}
		\Nilpotent(\lie{h}^\perp, H) = \Nilpotent((\lie{h}_\nc)^\perp, H_\nc) = \Nilpotent((\lie{h}')^\perp, H_\nc).
	\end{align*}
\end{lemma}

\begin{proof}
	Since $H/H_\nc$ is compact, we have $\Nilpotent(\lie{g}, H) = \Nilpotent(\lie{g}, H_\nc)$.
	Hence we may replace $H$ with $H_\nc$ in the assertion.
	
	Set $W\coloneq \lie{h}_\nc^\perp \cap \lie{h}^{H_{\nc}}$.
	Let $p\colon \lie{g} \rightarrow W$ be the orthogonal projection.
	Then $p$ is $H_{\nc}$-equivariant, and hence we have
	\begin{align*}
		p(\Nilpotent(\lie{h}_\nc^\perp, H_{\nc})) \subset \Nilpotent(W, H_{\nc}) = \set{0}.
	\end{align*}
	This shows that $\Nilpotent(\lie{h}_\nc^\perp, H_{\nc}) \subset \lie{h}^\perp$, which implies the first equality.
	The second equality follows from the same argument as above.
\end{proof}

Set $\lie{h''} \coloneq \lie{h} + \lie{k}_{H'}$ and $H'' \coloneq H K_{H'}$.
Then $K_{H'}$ is a maximal compact subgroup of $H''$ and $H''/H_\nc$ is compact.

\begin{lemma} \label{lem:NilpotentAndInvariant}
	One has $((\lie{h''})^\perp)^{H_\nc} \cap \Nilpotent(\lie{g}, G) = \set{0}$.
\end{lemma}

\begin{proof}
	Since $\lie{h'} = \lie{h''} + (\lie{h'})^{-\theta}$, we have
	\begin{align*}
		((\lie{h''})^\perp)^{H_\nc} \subset (\lie{h''})^\perp \cap \lie{h'} \subset (\lie{h'})^{-\theta}.
	\end{align*}
	Hence any element of $((\lie{h''})^\perp)^{H_\nc}$ is semisimple and is therefore not nilpotent unless it is zero.
	This proves the assertion.
\end{proof}

\begin{corollary} \label{cor:NilpotentAndPerp}
	Let $\orbit{O}$ be a nilpotent $G$-orbit in $\lie{g}$.
	Then $\overline{\orbit{O}} \cap \Nilpotent((\lie{h''})^\perp, H'') = \set{0}$ if and only if $\overline{\orbit{O}} \cap (\lie{h''})^\perp = \set{0}$.
\end{corollary}

\begin{proof}
	The assertion follows from Proposition \ref{prop:IntersectionW} and Lemma \ref{lem:NilpotentAndInvariant}.
\end{proof}

The following lemma was proved in \cite[Proposition 22]{Ki24} in a more general setting.

\begin{lemma} \label{lem:DiscreteDecomposabilityAndAdmissibilityWavefront}
	Let $\orbit{O}$ be a nilpotent $G$-orbit in $\lie{g}$.
	Then the following conditions are equivalent.
	\begin{enumparen}
		\item $\Nilpotent(\lie{h}^\perp, H) \cap \overline{\orbit{O}} = \set{0}$.
		\item $\Nilpotent(\lie{h}_\nc^\perp, H_\nc) \cap \overline{\orbit{O}} = \set{0}$.
		\item $(\lie{h}'')^\perp \cap \overline{\orbit{O}} = \set{0}$.
		\item $\lie{k}_{H'}^\perp \cap \overline{\orbit{O}} = \set{0}$.
		\item $(\lie{h}')^\perp \cap \overline{\orbit{O}} = \set{0}$.
	\end{enumparen}
\end{lemma}

\begin{proof}
	(1) $\Leftrightarrow$ (2) follows from Lemma \ref{lem:NilpotentAndNonCompact}.
	The equivalence of (3), (4) and (5) has been proved in Proposition \ref{prop:AdmissibilityCriterionWavefront}.
	Applying Lemma \ref{lem:NilpotentAndNonCompact} to $H''$, we have
	\begin{align*}
		\Nilpotent(\lie{h}_\nc^\perp, H_\nc) = \Nilpotent((\lie{h''})^\perp, H'').
	\end{align*}
	By Corollary \ref{cor:NilpotentAndPerp}, we have
	\begin{align*}
		\overline{\orbit{O}} \cap \Nilpotent(\lie{h}_\nc^\perp, H_\nc) = \set{0}
		&\Leftrightarrow \overline{\orbit{O}} \cap \Nilpotent((\lie{h''})^\perp, H'') = \set{0} \\
		&\Leftrightarrow \overline{\orbit{O}} \cap (\lie{h''})^\perp = \set{0}.
	\end{align*}
	This shows that (2) $\Leftrightarrow$ (3).
	We have shown the lemma.
\end{proof}

\begin{theorem} \label{thm:DiscretelyDecomposableAndAdmissibility}
	Let $V$ be an irreducible $(\lie{g}, K)$-module.
	Then the following conditions are equivalent.
	\begin{enumparen}
		\item $V|_{\lie{h}, K_H}$ is discretely decomposable.
		\item $V$ is $\lie{h}''$-admissible.
		\item $V$ is $K_{H'}$-admissible.
		\item $V$ is $\lie{h}'$-admissible.
	\end{enumparen}
\end{theorem}

\begin{proof}
	The assertion follows from Theorems \ref{thm:MainTheorem} and \ref{thm:AdmissibilityCriterion} and Lemma \ref{lem:DiscreteDecomposabilityAndAdmissibilityWavefront}.
\end{proof}

Theorem \ref{thm:DiscretelyDecomposableAndAdmissibility} asserts that there is no essential difference between admissibility and discrete decomposability.
For example, the following proposition gives a situation in which these two conditions are equivalent.

\begin{proposition} \label{prop:AdmissibilityAndDiscreteDecomposability}
	Let $V$ be an irreducible $(\lie{g}, K)$-module.
	Assume that $(\lie{h}^\perp)^{H_\nc} \cap \WF(V) = \set{0}$.
	Then $V$ is $\lie{h}$-admissible if and only if $V|_{\lie{h}, K_H}$ is discretely decomposable.
\end{proposition}

\begin{remark}
	Clearly, if $(\lie{h}^\perp)^{H_\nc} \cap \Nilpotent(\lie{g}, G) = \set{0}$, then the assumption holds.
	For example, $\lie{h} = \lie{h}''$ implies $(\lie{h}^\perp)^{H_\nc} \cap \Nilpotent(\lie{g}, G) = \set{0}$ as we have seen in Lemma \ref{lem:NilpotentAndInvariant}.
	We have given a necessary and sufficient condition for $(\lie{h}^\perp)^{H_\nc} \cap \Nilpotent(\lie{g}, G) = \set{0}$ in \cite[Proposition 22]{Ki24}.
\end{remark}

\begin{proof}
	By assumption and Proposition \ref{prop:IntersectionW}, we have
	\begin{align*}
		\WF(V) \cap \lie{h}^\perp = \set{0} \Leftrightarrow \WF(V) \cap \Nilpotent(\lie{h}^\perp, H) = \set{0}.
	\end{align*}
	Hence the assertion follows from Theorems \ref{thm:MainTheorem} and \ref{thm:AdmissibilityCriterion}.
\end{proof}

As in the proof, the conclusion of Proposition \ref{prop:AdmissibilityAndDiscreteDecomposability} holds under the condition
\begin{align*}
	\WF(V) \cap \lie{h}^\perp = \set{0} \Leftrightarrow \WF(V) \cap \Nilpotent(\lie{h}^\perp, H) = \set{0}.
\end{align*}
We have given a sufficient condition for this equivalence in \cite[Theorem 23]{Ki24}.

\subsection{Alternative proof of Fact \ref{fact:Implication} \texorpdfstring{(1) $\Rightarrow$ (2)}{(1) => (2)}}

In \cite{Ki24}, we proved the implication (1) $\Rightarrow$ (2) in Fact \ref{fact:Implication} by using several tools, e.g., heat kernels, wave front sets and translation functors.
Motivated by the insight obtained from Theorem \ref{thm:DiscretelyDecomposableAndAdmissibility}, we give an alternative proof of the implication (1) $\Rightarrow$ (2) in Fact \ref{fact:Implication}.
Retain the notation from the previous subsection.

We use the following fact, which was proved by T.\ Kobayashi in \cite[Theorem 2.9]{Ko98_discrete_decomposable_2}.
The original form of the fact is stated in terms of asymptotic $K$-support.
We shall state the fact in terms of the wave front set.
See \cite[Corollary 9]{Ki24} for the equivalence of the two conditions.

\begin{fact} \label{fact:AdmissibilityCriterionCompact}
	Let $V$ be an irreducible $(\lie{g}, K)$-module.
	If $\WF(V) \cap \lie{k}_{H}^\perp = \set{0}$, then $V$ is $K_{H}$-admissible and, in particular, $\lie{h}$-admissible.
\end{fact}

Although Fact \ref{fact:AdmissibilityCriterionCompact} is a special case of Theorem \ref{thm:AdmissibilityCriterion}, we have stated it separately to make explicit the results needed to prove the following theorem.

\begin{theorem}
	Let $V$ be an irreducible $(\lie{g}, K)$-module.
	If $\WF(V) \cap \Nilpotent(\lie{h}^\perp, H) = \set{0}$, then $V|_{\lie{h}, K_H}$ is discretely decomposable.
\end{theorem}

\begin{proof}
	Assume that $\WF(V) \cap \Nilpotent(\lie{h}^\perp, H) = \set{0}$.
	By Lemma \ref{lem:DiscreteDecomposabilityAndAdmissibilityWavefront}, we have $\WF(V) \cap (\lie{k}_{H'})^\perp = \set{0}$.
	Hence $V$ is $\lie{h}''$-admissible by Fact \ref{fact:AdmissibilityCriterionCompact}.
	Take an $(\lie{h}'', K_{H'})$-module filtration $V = \bigcup_{i=0}^\infty V_i$ such that each $V_i$ has finite length.
	From the definition of $\lie{h}''$ and $\lie{h}_\nc$, there exists a compact ideal $\lie{k}''$ of $\lie{h}''$ such that
	\begin{align*}
		\lie{h}_\nc \subset \lie{h} \subset \lie{h}_{\nc} \oplus \lie{k}'' = \lie{h}''.
	\end{align*}
	This implies that $V_i$ has finite length as an $(\lie{h}, K_{H})$-module for each $i$.
	Hence $V|_{\lie{h}, K_H}$ is discretely decomposable.
\end{proof}